% !PRINTED 329
\weberpage{329}

Lassen wir aber nun $x$ sich der Wurzel $\alpha$ annähern, so
wird $X$ sich derselben Wurzel nähern, und (20) zeigt, dass auch
$\xi$ derselben Grenze zustrebt. Es folgt also, dass, wenn $x$ hinlänglich
nahe an einer Wurzel liegt, diese nächste Wurzel in
dem Intervall $(\xi, X)$ liegt, in dem keine zweite Wurzel enthalten
ist, und dass dies Intervall sich mehr und mehr um $\alpha$ schliesst.

Wir nehmen als Beispiel (mit \gesperrt{Laguerre}) die Kugelfunctionen,
von denen wir schon früher (\S.\,87) nachgewiesen haben,
dass sie nur reelle Wurzeln haben, die alle zwischen $-1$ und
$+1$ liegen. Die Function
\[
P_n(x) = \frac{1\,.\,3 \ldots (2\,n-1)}{1\,.\,2 \ldots n} \left( x^n - \frac{n(n-1)}{2(2\,n-1)}\,x^{n-2} + \cdots \right)
\]
hat die Eigenschaft, dass
\[
P_n(1) = 1, \qquad P_n'(1) = \frac{n(n+1)}{2},
\]
\[
P_n''(1) = \frac{n(n-1)(n+1)(n+2)}{8}
\]
ist. Man beweist dies leicht durch vollständige Induction, indem
man in den ersten Fällen $n = 2, 3 \ldots$ die Formeln direct
nachweist, und dann aus der Recursionsformel
\[
(1 - x^2)\,P_n'(x) + n\,x\,P_n(x) - n\,P_{n-1}(x) = 0
\]
durch zweimalige Ableitung die Richtigkeit für den Index $n$ nachweist,
falls sie für den Index $n-1$ vorausgesetzt wird.

Wenn wir dann in (19) und (20) $x = 1$ setzen, so erhalten
wir zwei Grenzen, zwischen denen die der 1 am nächsten kommende
Wurzel liegt, nämlich
\[
X = 1 - \frac{2}{n + 1 + (n-1)\sqrt{\dfrac{n(n+1)}{2}}},
\]
\[
\xi = 1 - \frac{2}{n + 1 + \sqrt{\dfrac{n(n+1)}{2}}}.
\]
Für $n = 7$ z. B. ergiebt sich
\[
X = 0{,}94967 \ldots
\]
\[
\xi = 0{,}84952 \ldots
\]
